%!TEX root=thesis.tex
Cardiac Magnetic Resonance (CMR) imaging is widely used for quantitative assessment of cardiac anatomy and function. In automated cardiac analysis, accurate segmentation of the right ventricle (RV), myocardium (MYO), and left ventricle (LV) is an important prerequisite for deriving functional measurements such as end-diastolic volume, end-systolic volume, myocardial mass, and ejection fraction. Although U-Net-style convolutional neural networks remain strong foundations for medical image segmentation, practical deployment also requires attention to model size, inference cost, and reproducible validation. This thesis investigates Nano-Mamba U-Net, a lightweight Mamba-inspired 3D U-Net architecture for cardiac MRI segmentation. The model keeps the encoder-decoder structure and skip connections of U-Net, while inserting a compact sequence-processing bottleneck at the compressed latent representation. This bottleneck serializes 3D feature maps into a one-dimensional sequence, applies a lightweight gated sequence module inspired by Mamba-style state space modeling, and reshapes the output back into volumetric form. The implementation is therefore described as Mamba-inspired rather than as a full reproduction of the original hardware-aware selective scan algorithm.

The proposed model was evaluated on the Automated Cardiac Diagnosis Challenge (ACDC) dataset using a deterministic patient-level validation protocol. The final experiment used 80 training patients and 20 held-out validation patients, corresponding to 160 training cases and 40 validation cases. All compared models were trained under the same preprocessing and validation procedure, and the best checkpoint for each model was selected by validation mean Dice. Nano-Mamba U-Net achieved a validation mean Dice Similarity Coefficient (DSC) of 84.78\% with 1.46 million trainable parameters. It improved over the 3D U-Net baseline, which achieved 80.83\% mean DSC with 4.81 million parameters. However, SegResNet16 achieved the highest mean DSC of 86.70\%, and the No-Mamba ablation achieved 85.64\%. These results show that Nano-Mamba U-Net should not be claimed as the most accurate model by Dice. Instead, its main contribution is an accuracy-efficiency trade-off: it provides competitive held-out validation performance while using the smallest parameter count among the evaluated models. The thesis concludes that lightweight sequence-inspired bottlenecks are promising for cardiac MRI segmentation, but further work is required for official test-set evaluation, pathology-specific analysis, external validation, and full 4D cardiac motion modeling.

\vspace{1em}
\noindent
\textbf{Keywords:} Cardiac MRI Segmentation, Deep Learning, Mamba-Inspired Architecture, State Space Models, Accuracy-Efficiency Trade-off.
